\documentclass{article}
\usepackage[UTF8]{ctex}
\usepackage{amsmath}
\usepackage{amssymb}
\usepackage{geometry}
\geometry{a4paper, margin=2.5cm}

\title{任务空间控制中的坐标系记号说明}
\author{}
\date{}

\begin{document}

\maketitle

\section{问题：$X$是什么？}

在第432行提到：
\begin{quote}
设末端执行器轨迹由$(X(t), V_b(t))$指定，其中$X \in SE(3)$且$[V_b] = X^{-1}\dot{X}$，即twist $V_b$在末端执行器坐标系$\{b\}$中表示。
\end{quote}

\textbf{答案}：$X$是从\textbf{空间坐标系$\{s\}$}到\textbf{体坐标系$\{b\}$}（末端执行器坐标系）的变换矩阵，即$X = X_{sb}$。

\section{标准记号}

在现代机器人学中，通常使用以下记号：

\begin{itemize}
\item $X_{ab}$：从坐标系$\{a\}$到坐标系$\{b\}$的齐次变换矩阵
\item $\{s\}$：空间坐标系（space frame），通常是固定坐标系
\item $\{b\}$：体坐标系（body frame），这里指末端执行器坐标系
\end{itemize}

因此：
\begin{equation}
X = X_{sb} = \begin{bmatrix}
R_{sb} & p_{sb} \\
0 & 1
\end{bmatrix} \in SE(3),
\end{equation}
其中：
\begin{itemize}
\item $R_{sb} \in SO(3)$：从$\{s\}$到$\{b\}$的旋转矩阵
\item $p_{sb} \in \mathbb{R}^3$：从$\{s\}$到$\{b\}$的平移向量
\end{itemize}

\section{为什么使用$X$而不是$X_{sb}$？}

在任务空间控制的上下文中，为了\textbf{记号简洁}，通常省略下标，直接用$X$表示末端执行器相对于空间坐标系的配置。这是常见的约定，因为：

\begin{enumerate}
\item 在任务空间控制中，我们主要关心末端执行器的位姿
\item 空间坐标系$\{s\}$通常是固定的参考坐标系
\item 体坐标系$\{b\}$是末端执行器坐标系
\item 省略下标使公式更简洁易读
\end{enumerate}

\section{从上下文确认}

\subsection{第432行的上下文}

\begin{quote}
设末端执行器轨迹由$(X(t), V_b(t))$指定，其中$X \in SE(3)$且$[V_b] = X^{-1}\dot{X}$，即twist $V_b$在末端执行器坐标系$\{b\}$中表示。
\end{quote}

\textbf{关键信息}：
\begin{itemize}
\item $V_b$是体坐标系twist（body-frame twist）
\item $[V_b] = X^{-1}\dot{X}$是体坐标系twist的标准定义
\item 这明确表明$X$是从空间坐标系到体坐标系的变换
\end{itemize}

\subsection{第434行的确认}

\begin{quote}
正运动学是$X = T(\theta)$和$V_b = J_b(\theta)\dot{\theta}$。
\end{quote}

\textbf{说明}：
\begin{itemize}
\item $T(\theta)$是正运动学函数，计算末端执行器相对于基座（空间坐标系）的位姿
\item $J_b(\theta)$是体雅可比矩阵（body Jacobian），计算体坐标系twist
\item 这进一步确认$X$是相对于空间坐标系的配置
\end{itemize}

\subsection{第513行的明确说明}

在混合运动-力控制部分，有更明确的说明：
\begin{quote}
设$X(t) \in SE(3)$是块的坐标系$\{b\}$相对于空间坐标系$\{s\}$的配置。
\end{quote}

这\textbf{明确}说明了$X$是从$\{s\}$到$\{b\}$的变换。

\section{体坐标系Twist的定义}

体坐标系twist $V_b$的定义为：
\begin{equation}
[V_b] = X^{-1}\dot{X} = X_{sb}^{-1}\dot{X}_{sb}.
\end{equation}

\textbf{物理意义}：
\begin{itemize}
\item $V_b$表示末端执行器的瞬时运动（角速度和线速度）
\item 这些速度在末端执行器自身的坐标系$\{b\}$中表示
\item $X^{-1}\dot{X}$是从变换矩阵的时间导数计算体坐标系twist的标准方法
\end{itemize}

\section{与空间坐标系Twist的关系}

空间坐标系twist $V_s$定义为：
\begin{equation}
[V_s] = \dot{X}X^{-1} = \dot{X}_{sb}X_{sb}^{-1}.
\end{equation}

\textbf{关系}：
\begin{equation}
V_s = [Ad_X]V_b,
\end{equation}
其中$[Ad_X]$是伴随表示（adjoint representation），用于在不同坐标系之间转换twist。

\section{总结}

\begin{enumerate}
\item \textbf{$X$的含义}：$X = X_{sb}$，从空间坐标系$\{s\}$到体坐标系$\{b\}$（末端执行器坐标系）的变换矩阵

\item \textbf{记号约定}：在任务空间控制中，为了简洁，通常省略下标，直接用$X$表示

\item \textbf{体坐标系twist}：$V_b$是体坐标系twist，通过$[V_b] = X^{-1}\dot{X}$计算

\item \textbf{正运动学}：$X = T(\theta)$表示末端执行器相对于基座（空间坐标系）的位姿

\item \textbf{上下文确认}：从多个地方的上下文可以确认$X$就是$X_{sb}$
\end{enumerate}

\section{常见记号对照}

\begin{table}[h]
\centering
\begin{tabular}{|l|l|}
\hline
\textbf{完整记号} & \textbf{简化记号（本文档）} \\
\hline
$X_{sb}$ & $X$ \\
\hline
$V_b$（体坐标系twist） & $V_b$ \\
\hline
$V_s$（空间坐标系twist） & $V_s$ \\
\hline
$F_b$（体坐标系wrench） & $F_b$ \\
\hline
$F_s$（空间坐标系wrench） & $F_s$ \\
\hline
\end{tabular}
\caption{坐标系记号对照表}
\end{table}

\textbf{注意}：虽然$X$省略了下标，但$V_b$、$F_b$等仍然保留下标，因为它们明确表示在哪个坐标系中。

\end{document}

